\documentclass[11pt,a4paper]{article}
\usepackage[utf8]{inputenc}
\usepackage[british]{babel}
\usepackage{geometry}
\geometry{margin=1in}
\usepackage{graphicx}
\usepackage{booktabs}
\usepackage{longtable}
\usepackage{float}
\usepackage{amsmath}
\usepackage{hyperref}
\usepackage{algorithm}
\usepackage{algpseudocode}

\title{Risk Management \\ Group Work Project 1}
\author{Ge Zhang \\ Alberto Hernandez-Espinosa}
\date{\today}

\begin{document}

\maketitle



\section*{On AI Usage}

This research made use of artificial intelligence tools in three 
specific capacities: (i) to assist with literature search and verification 
of academic sources and datasets; (ii) to help refine and optimize algorithmic 
implementations during the structure learning phase; and (iii) for grammar, 
spelling, and style corrections, particularly for itemized content where precision 
in phrasing mattered.

All intellectual contributions, theoretical framing, empirical analysis, 
and final interpretations are entirely our own. The AI tools served only as 
supporting aids, not as co-authors, and we assume full responsibility for any errors or omissions.



\section{Problem Formulation}
The global crude oil market is among the most complex and closely watched commodity systems out 
there. Unlike your typical financial equities, crude oil’s spot price isn’t just about 
discounting expected future cash flows—it acts more like a vital macroeconomic pulse, 
shaping industrial costs, swaying inflation, and even nudging geopolitical decisions. 
So, trying to forecast crude oil prices isn’t really a straightforward time-series 
exercise. What’s needed is a broader framework, one that can tease apart 
how macroeconomic forces, supply-chain mechanics at the micro level, and 
unexpected geopolitical events all intertwine.

\subsection{Limitations of Traditional Econometric Approaches}

Looking back at how quantitative analysts and economists have approached 
this problem, they've tended to lean on familiar workhorses like ARIMA 
and GARCH models. There's no denying these are mathematically solid and 
they do a decent job picking up things like volatility clustering or mean reversion. 
But if we apply them to today's oil market — which is driven by so many interconnected 
variables — they start to show their weaknesses.

One big issue is that these models are fundamentally atheoretical. They're 
good at finding statistical correlations but they don't really tell you anything 
about causation. The GARCH model, for example, it might predict that volatility 
will carry on being volatile, but it can't explain \textit{why} a supply 
disruption in the Middle East ends up affecting refinery margins in the U.S. 
The causal pathway just isn't there.

Another problem is how they handle data at different frequencies and 
structural breaks. If we get macroeconomic data like GDP or CPI coming out monthly 
or quarterly, while futures prices are updating every second. Mixing those 
together is messy. And then you've got these sudden regime shifts, the 2014 shale 
revolution, or the demand collapse during COVID, which completely break the 
stationarity assumptions that these classical models depend on.

Plus, we can see that these models don't handle missing data well at all. 
And when we put aññ together global microeconomic data, especially from countries 
that aren't exactly transparent about their inventory levels, missing values are 
pretty much the norm.

\subsection{The Suitability of Probabilistic Graphical Models}

To get around these limitations, the proposal is a Probabilistic Graphical Models (PGMs), 
specifically Belief Networks (Bayesian Networks). PGMs give us a elegant way to combine 
graph theory with probability theory, which lets us represent complex, 
high-dimensional joint distributions in a compact form.

A Belief Network is basically a Directed Acyclic Graph (DAG). The nodes are 
 random variables, as OPEC production quotas, interest rates from the Fed, 
 or the Brent-WTI spread, and the directed edges show conditional dependencies 
 and causal influences. What makes this architecture so powerful is the local 
 Markov property: once you condition on its parents, a node becomes independent 
 of all its non-descendants. This drastically cuts down the computational complexity, 
 turning what would be an intractable joint distribution into a product of much 
 simpler conditional probability distributions.

\subsection{Advantages of the Proposed Methodology}

There are several concrete advantages to using Belief Networks for forecasting 
crude oil prices, especially when we look at it through the three lenses of 
macroeconomic, microeconomic, and financial specialisation.

\subsubsection{Advantages of the Proposed Methodology}

\textbf{Interpretability:} Unlike deep neural networks or other black-box approaches, 
PGMs give us transparent causal pathways that you can actually trace and understand. 
Through structure learning, the model can spot hidden dependencies in the data 
that might either confirm or challenge existing economic theories. 
For someone managing a portfolio or making policy decisions, being able to 
visualise exactly how a shock to the US Dollar Index flows through to oil demand 
in emerging markets, that's invaluable for running scenarios and stress tests.

\textbf{handling of uncertainty and missing data}. Because Belief Networks are 
probabilistic by nature, they keep working even when some variables aren't observed. 
Say weekly crude stock data is late or just not available, the network can still 
marginalise over that uncertainty and give you a probability distribution for the 
spot price. That's way more useful than a single point estimate that might be completely off.

\textbf{Blend expert knowledge with data-driven learning: } 
The algorithm figure out the structure from the data, but  
also impose constraints based on what we know about the world. 
For example, it might be possible to 
specify that geopolitical risk can influence supply, but actually the supply can't 
 change past geopolitical events. That kind of guidance helps the model avoid 
picking up spurious correlations that don't make economic sense.

\textbf{Regime-switching angle}. Combining 
PGMs with latent variable models like Hidden Markov Models, this framework 
could be adapted to a different market conditions on the fly, whether we're in contango 
or backwardation, or during high-volatility crisis periods versus calmer growth phases. 
Energy markets are notoriously non-linear, and this gives us a way to 
capture those changing dynamics.

So overall, the idea is to build something that's not just statistically 
sophisticated but also causally informed and robust to the kind of messiness 
you inevitably run into with real-world data.

To sum up, the core aim here is to move beyond purely statistical time-series 
forecasting, building a framework that's causally aware and probabilistically robust. 
By learning both the structure and parameters of a Belief Network from a 
dataset that spans macroeconomic, microeconomic, and financial domains.










\section{Data Identification, Sources and Perspectives}

We broked down the data  
into three distinct angles: the macroeconomic and geopolitical side, 
the microeconomic side, and the financial side. Each one brings 
something different to the table.

\subsection{Macroeconomic and Geopolitical Perspective}

This lens is all about the big picture: sovereign policies, 
global systemic drivers, that sort of thing. We're looking at variables 
like the Consumer Price Index (CPI) and the Federal Funds Rate, which 
provide a sense of the broader inflationary environment and the cost of capital. 
Those are the kinds of things that shape global aggregate demand. Industrial 
Production is another key one; it's a pretty direct measure of real 
economic activity and how much energy the manufacturing sector is consuming. 
We've also included the US Dollar Index, and that's because crude oil is priced 
in dollars globally. So when the dollar's purchasing power shifts, it 
feeds back into nominal oil prices, usually in the opposite direction, 
and changes how demand behaves in emerging markets.

\subsection{Microeconomic Perspective}

Trying to gethigh-frequency microeconomic data, like individual 
refinery margins or real-time rig counts is a challenge by itself. 
Reporting lags, proprietary restrictions, it's a whole thing. But even without 
that granular level of detail, we can still capture the physical market balance 
through the WTI spot price and the futures curves. The relationship between 
commercial inventory levels and extraction costs tends to shape short-term 
physical supply responses, and that ultimately shows up in the volatility of 
spot and near-month futures contracts.

\subsection{Financial Perspective}

We're using front-month WTI and Brent crude futures 
contracts to get at market sentiment, where speculative positions are 
sitting, and the geographical arbitrage play, that's the Brent-WTI spread. 
These instruments reflect how the market prices risk at the margin, 
and they tend to react pretty sharply to macroeconomic news and geopolitical developments.





\section{Data Dictionary and Summary Table}

We pulled the raw data from two main sources. For the high-frequency market stuff, 
we went with Yahoo Finance, because it's reliable and gives us good coverage. 
For the macroeconomic indicators, we used the Federal Reserve Economic Data (FRED) 
repository, which is pretty much the gold standard for that kind of thing. 

Now, one challenge was getting everything to line up properly. The market data 
comes in at daily frequency, but the macro stuff is monthly. So we aligned all 
the daily observations to a monthly frequency using the Month-End convention 
(that's the \texttt{ME} approach). This way everything syncs up nicely with 
the macroeconomic releases. This resulted in a balanced panel of 199 monthly 
observations, covering January 2010 right through to the present. 

Table \ref{tab:data_summary} below gives a full dictionary of all the variables 
we've integrated into the unified dataset.


\begin{table}[H]
\centering
\caption{Data Dictionary and Summary of Key Variables}
\label{tab:data_summary}
\renewcommand{\arraystretch}{1.3}
\resizebox{\textwidth}{!}{%
\begin{tabular}{@{}lllll@{}}
\toprule
\textbf{Variable (Ticker)} & \textbf{Description} & \textbf{Native Freq.} & \textbf{Source} & \textbf{Date Range} \\
\midrule
\texttt{WTI\_CL} & WTI Crude Oil Futures (Front Month) & Daily & Yahoo Finance & 2010-01 to Present \\
\texttt{Brent\_BZ} & Brent Crude Oil Futures (Front Month) & Daily & Yahoo Finance & 2010-01 to Present \\
\texttt{USD\_Idx} & US Dollar Index (Basket of Currencies) & Daily & Yahoo Finance & 2010-01 to Present \\
\texttt{CPI} & Consumer Price Index for All Urban Consumers & Monthly & FRED (CPIAUCSL) & 2010-01 to Present \\
\texttt{Fed\_Funds} & Effective Federal Funds Rate & Monthly & FRED (FEDFUNDS) & 2010-01 to Present \\
\texttt{Ind\_Prod} & Industrial Production Index & Monthly & FRED (INDPRO) & 2010-01 to Present \\
\texttt{WTI\_Spot} & Spot Price of WTI Crude Oil & Daily & FRED (DCOILWTICO) & 2010-01 to Present \\
\bottomrule
\end{tabular}%
}
\end{table}





\section{Data Preprocessing and Cleaning}

So we followed a sterilisation protocol of athree-stage cleaning process:
\begin{enumerate}
    \item Extreme Outlier Identification and Clipping
    \item Bad Data and Duplicated Indices
    \item Missing Values and Temporal Interpolation
\end{enumerate}

\subsection{Extreme Outlier Identification and Clipping}

We went with the Interquartile Range (IQR) method. The bounds were 
set at $Q1 - 3 \times IQR$ and $Q3 + 3 \times IQR$. 
We used factor of 3 used for standard boxplots because commodity returns 
are inherently leptokurtic, they have these fat tails where extreme 
price movements, while statistically rare, are actually 
economically meaningful. Considering for example the 2008 financial crisis 
or the COVID demand shock in 2020; those weren't outliers in the sense of 
"bad data," they were real events.

When we ran the preprocessing pipeline, only two extreme values 
got flagged and clipped across the whole dataset. That's a pretty low 
number, which tells us the FRED and Yahoo Finance feeds are good quality. 
The values we clipped were most likely flash-crash anomalies or 
erroneous ticks that the exchanges later corrected, not genuine structural market events.

\subsection{Bad Data and Duplicated Indices}

W aligned everything to the last business day of each month, then the chances 
of duplicated chronological indices were already pretty low. 
We confirm with zero duplicated indices in both the raw and aggregated dataframes, 
so we ended up with a strictly monotonic time-series progression. No issues there.

\subsection{Missing Values and Temporal Interpolation}

We used time-based interpolation for to go with mean or median imputation. 
By this approach we aimed to destroy the temporal autocorrelation structure,
which is kind of a big deal for financial time series. Time-based interpolation, 
on the other hand, respects the chronological distance between known data points. For 
any remaining missing values at the very beginning or end of the dataset, 
we used forward-filling and backward-filling. After all this, we had a complete, 
balanced panel dataset with no null values, ready for exploratory analysis and probabilistic modelling.

\begin{figure}[htbp]
\centering
\includegraphics[width=\textwidth]{figures/cleaned_wti_price.png}
\caption{Sterilised WTI Crude Oil Futures Price (2010--Present). The blue line represents the robust, outlier-clipped series, demonstrating the preservation of major structural breaks such as the 2014 shale glut and the 2020 pandemic crash.}
\label{fig:cleaned_wti}
\end{figure}





\section{Exploratory Data Analysis}

Once we had the sterilised monthly dataset put together, we dove into 
a thorough Exploratory Data Analysis (EDA) in order to the underlying 
statistical properties of the crude oil market. We broke it down into 
distributional, temporal, and multivariate dimensions, and by the end of it, 
we'd identified the key stylised facts that tend to characterise energy commodities.

\subsection{Distributional and Temporal Analysis}

Figure \ref{fig:returns_dist} a non-Gaussian. The monthly percentage returns 
have this pronounced leptokurtic shape. This is, there's a sharp peak around the mean, 
but then these heavy, asymmetric tails on either side. So oil prices can 
cruise along with relative stability for months at a time, but they're also 
prone to these sudden, violent shocks in either direction.

\begin{figure}[htbp]
\centering
\includegraphics[width=0.8\textwidth]{figures/returns_distribution.png}
\caption{Distribution of Monthly WTI Percentage Returns. The pronounced kurtosis and fat tails demonstrate the failure of normal distribution assumptions in energy markets.}
\label{fig:returns_dist}
\end{figure}

The time-series plot of the sterilised WTI price (Figure \ref{fig:cleaned_wti}) 
tells an interesting story: several distinct market regimes playing out. 
From 2010 to 2014, prices were consistently high, fuelled by strong demand 
from emerging markets. Then came the structural break in late 2014, that's when 
the US shale revolution really kicked in, leading to a supply glut. 
And looking at early 2020, that's where things get really wild. 
Unprecedented demand destruction met storage capacity constraints, and futures 
prices briefly went negative. Something not very common.

\subsection{Multivariate Dependencies}

To get a sense of how the macroeconomic, microeconomic, and financial variables 
interact with each other, we generated a correlation matrix of the monthly 
percentage returns.

Figure \ref{fig:corr_heatmap} confirms what economic theory would predict. The 
most obvious relationship is the persistent negative correlation between 
the US Dollar Index and both WTI and Brent returns. Makes sense when 
considering oil's priced in dollars, so when the dollar strengthens, oil becomes 
more expensive for anyone holding foreign currency. That tends to dampen 
demand and push prices down. On the flip side, industrial production and 
inflation metrics generally show positive correlations with energy prices, 
though there's usually a lag. That reflects the derived demand for crude 

\begin{figure}[htbp]
\centering
\includegraphics[width=0.8\textwidth]{figures/correlation_heatmap.png}
\caption{Correlation Heatmap of Monthly Returns. Notable inverse relationships are observed between the US Dollar Index and crude oil returns, validating the financial transmission mechanism of currency fluctuations.}
\label{fig:corr_heatmap}
\end{figure}



\subsection{Stylised Facts of Oil Prices}

Based on everything we saw in the EDA, we can pin down some stylised facts 
about crude oil prices. Any forecasting model worth its salt, including our 
proposed Belief Network — needs to be able to capture these:

\begin{enumerate}
    \item \textbf{Volatility Clustering and Structural Breaks:} Unlike standard 
    equity markets, oil prices are highly susceptible to exogenous supply shocks 
    and geopolitical events. This results in distinct periods of extreme 
    volatility clustering, separated by structural breaks where the underlying 
    data-generating process fundamentally shifts (e.g., the transition 
    from a supply-constrained market to a demand-constrained market).
    \item \textbf{Non-Normal Return Distributions:} Oil returns are distinctly 
    non-Gaussian. They exhibit high kurtosis (fat tails) and skewness, 
    meaning that extreme events (e.g., a $>20\%$ monthly move) occur far 
    more frequently than a normal distribution would predict.
    \item \textbf{Weak Autocorrelation in Returns, Strong Persistence in Volatility:} 
    The raw monthly returns show very low linear autocorrelation, 
    suggesting that the market is highly efficient and past price 
    directions do not easily predict future directions. However, the squared 
    returns (a proxy for volatility) exhibit strong positive autocorrelation. 
    This confirms that while the direction of the price change is unpredictable, 
    the magnitude of the change is highly persistent.
    \item \textbf{Macroeconomic and Financial Sensitivity:} Oil prices do not 
    exist in a vacuum. They are deeply sensitive to the cost of capital 
    (interest rates), global industrial activity, and currency fluctuations 
    (USD strength). The causal propagation of these variables is non-linear 
    and regime-dependent, necessitating a modelling framework capable of 
    handling complex conditional dependencies.
\end{enumerate}















\section{Probabilistic Graphical Models}

Looking at everything through those three lenses we've been using macroeconomic, 
microeconomic, and financial. The goal here is to lay down the mathematical 
and conceptual foundations we need to model the complex, interdependent 
dynamics of the global crude oil market.

\subsection{Macroeconomic and Geopolitical Perspective: Foundations of Probabilistic Graphical Models}

Looking at the global oil market from a macroeconomic and geopolitical angle, 
it's really just a giant web of interdependent variables. Sovereign policies, 
aggregate demand, geopolitical shocks, they all propagate through these intricate 
causal pathways. Probabilistic Graphical Models (PGMs) give us a way to 
represent that kind of system mathematically. They combine graph theory with 
probability theory, letting us compactly represent a joint probability 
distribution over a high-dimensional set of random variables. It map 
variables to nodes, probabilistic dependencies to edges for visualising and 
compute how macroeconomic shocks cascade through the system.

PGMs generally fall into two broad families: Belief Networks and Markov Networks.

\begin{itemize}
    \item \textbf{Belief Networks (Bayesian Networks):} These are directed acyclic graphs (DAGs). 
    The directed edges represent explicit causal influence or temporal precedence. 
    Each node comes with a conditional probability distribution that quantifies 
    the effect of its parent nodes. For oil markets, Belief Networks are particularly 
    well suited to modelling asymmetric causal chains. 
    Ths can be said: a geopolitical conflict (parent) disrupts supply chains (child), 
    which then drives up inflation (grandchild). The acyclic nature enforces a 
    logical flow of time and causality, which makes sense for these kinds of problems.
    \item \textbf{Markov Networks (Markov Random Fields):} These use undirected 
    graphs and are parameterised by potential functions over cliques of nodes, 
    rather than conditional probabilities. They're really effective for systems 
    with simultaneous, bidirectional feedback loops, situations where causal 
    directionality isn't clear. For example, the mutual influence between global 
    industrial production and aggregate energy consumption often happens 
    simultaneously within a given fiscal quarter. Undirected representations are 
    better at capturing those equilibrium states.
\end{itemize}

For forecasting crude oil prices, the goal is primarily leveraging Belief Networks. Being 
able to encode directional causal influence and perform both predictive (forward) 
and diagnostic (backward) reasoning is pretty essential for macroeconomic scenario analysis.

\subsection{Microeconomic Perspective: Parameter Learning versus Structure Learning}

From microeconomic lens, its possible to look at individual actors, refinery utilisation 
rates, inventory drawdowns, upstream capital expenditure, that sort of thing. 
To capture these dynamics accurately within a PGM, that we need to deal with in 
two distinct phases of model construction: parameter learning and structure learning.

\textbf{Parameter Learning} assumes you already know the structural topology 
of the graphical model, the DAG, maybe from established microeconomic theory or
 expert consensus. Then it's just a matter of estimating the numerical parameters 
 of the conditional probability distributions from the observed data. 
 This is done, typically with Maximum Likelihood Estimation (MLE) or Bayesian 
 estimation techniques. Say economic theory tells that commercial inventory 
 levels directly influence short-term spot price volatility. Parameter learning 
 will quantify the exact magnitude and variance of that relationship 
 based on historical microeconomic data.

\textbf{Structure Learning}, This is about discovering the underlying graphical 
topology directly from the data. A hard computational problem due the space of 
possible DAGs grows super-exponentially with the number of variables. Structure 
learning algorithms generally come in two flavours: constraint-based methods 
(which use statistical conditional independence tests to eliminate edges) 
and score-based methods (which search the graph space to optimise a goodness-of-fit 
metric penalised for complexity, like the Bayesian Information Criterion). 
For the crude oil market, structure learning is incredibly valuable. 
It lets the model autonomously uncover hidden, non-linear microeconomic 
dependencies that human analysts might miss, reducing reliance on potentially 
biased expert input.

\subsection{Financial Specialist Perspective: Markov Chains and Markov Blankets}

From a financial specialist's perspective, crude oil is volatile asset class. 
Modelling its temporal evolution and feature space requires specific probabilistic 
tools, namely Markov Chains and the concept of the Markov Blanket.

\textbf{Markov Chains and Hidden Markov Models (HMMs):} A Markov Chain is 
a stochastic model where the probability of each event depends only on the 
state attained in the previous event. That's the memoryless property. 
Now, raw oil prices don't strictly follow a simple Markov Chain because 
there are long-term memory effects. But the \textit{regimes} of the oil 
market often do follow one. Hidden Markov Models extend this by positing 
that the observed financial time series (say, daily returns) is generated 
by an underlying, unobservable Markov process representing distinct market 
states. So an HMM can dynamically classify the market into latent regimes, 
contango/low-volatility, backwardation/high-volatility, that kind of thing. 
This lets the forecasting model adapt its parameters dynamically based on 
the current financial regime.

\textbf{The Markov Blanket:} In a Belief Network, the Markov Blanket of 
a node $X$ is the minimal set of variables that completely shields $X$ 
from the rest of the network. Specifically, it consists of $X$'s parents, 
its children, and the other parents of its children (co-parents or spouses). 
The key property is that, conditional on this set, $X$ is conditionally 
independent of all other nodes. For a financial specialist building a 
forecasting model for the WTI spot price, identifying the Markov Blanket 
of the price node is essentially optimal feature selection. It isolates 
the exact subset of macroeconomic, microeconomic, and financial indicators 
that have genuine predictive power, while discarding redundant variables 
that just add noise to the trading algorithm.








\section{Structure Learning: The Inferred Causality Algorithm}

The structure of the Belief Network has to be learned from our sterilised 
dataset, and that's where the Inferred Causality algorithm comes in. 
It's a constraint-based approach, rigorous, systematic, designed to uncover 
the causal DAG by methodically evaluating conditional independencies.

The algorithm works through three sequential phases, and it leans 
heavily on that Markov Blanket concept. First, it identifies the Markov 
Blanket for every variable, which drastically cuts down the dimensionality 
of the search space. Second, it isolates the direct neighbours 
(parents and children) by filtering out the spouses. Finally, it orients 
the undirected edges to establish causal directionality, with special 
attention paid to colliders, those v-structures that tell us a lot about causality.

The formal pseudocode is laid out in Algorithm \ref{alg:inferred_causality} below.

\begin{algorithm}[H]
\caption{Inferred Causality Structure Learning}
\label{alg:inferred_causality}
\begin{algorithmic}[1]
\Require A dataset $\mathcal{D}$ containing $N$ variables $\mathbf{V} = \{X_1, X_2, \dots, X_N\}$, and a significance level $\alpha$ for conditional independence tests.
\Ensure A Directed Acyclic Graph (DAG) $\mathcal{G} = (\mathbf{V}, E)$.

\State \textbf{Phase 1: Markov Blanket Discovery}
\For{each variable $X_i \in \mathbf{V}$}
    \State Identify the Markov Blanket $\mathcal{MB}(X_i)$ using a forward-backward selection algorithm (e.g., Grow-Shrink or IAMB).
    \State Enforce symmetry: If $X_j \in \mathcal{MB}(X_i)$ but $X_i \notin \mathcal{MB}(X_j)$, remove the asymmetric dependency based on statistical reliability.
\EndFor

\State \textbf{Phase 2: Neighbourhood Identification (Undirected Skeleton)}
\For{each variable $X_i \in \mathbf{V}$}
    \State Initialise the set of Parents and Children $\mathcal{PC}(X_i) = \mathcal{MB}(X_i)$.
    \For{each $X_j \in \mathcal{PC}(X_i)$}
        \State Search for a separating set $\mathbf{S} \subseteq \mathcal{MB}(X_i) \setminus \{X_j\}$ such that $X_i \perp X_j \mid \mathbf{S}$.
        \If{such a separating set $\mathbf{S}$ is found}
            \State Remove $X_j$ from $\mathcal{PC}(X_i)$ (identify $X_j$ as a spouse/co-parent).
            \State Record $\mathbf{S}$ in the separation set table $\mathbf{Sep}(X_i, X_j)$.
        \EndIf
    \EndFor
\EndFor
\State Construct the undirected skeleton graph $\mathcal{G}_{skel}$ using the validated $\mathcal{PC}$ sets.

\State \textbf{Phase 3: Edge Orientation (V-Structures and Meek's Rules)}
\For{each unshielded triple $X_i - X_k - X_j$ in $\mathcal{G}_{skel}$}
    \If{$X_k \notin \mathbf{Sep}(X_i, X_j)$}
        \State Orient as a collider (v-structure): $X_i \rightarrow X_k \leftarrow X_j$.
    \EndIf
\EndFor
\State Apply Meek's propagation rules to orient remaining edges without creating directed cycles or new v-structures.
\State \Return The final oriented DAG $\mathcal{G}$.
\end{algorithmic}
\end{algorithm}




\bibliographystyle{plain}
\bibliography{references} % Add bib file later if needed

\end{document}